% ============================================================================
% Section 7.4 — Area of Composite Shapes
% CCSS 7.G.B.6
% ============================================================================
\section{Area of Composite Shapes}

\begin{stepGoal}
\begin{itemize}
    \item Break a composite shape into simpler pieces
    \item Add or subtract areas to find the total area
\end{itemize}
\end{stepGoal}

\begin{stepVocabBox}
    \vocabItem{Composite Shape}{A figure made by joining two or more simple shapes (rectangles, triangles, circles, etc.).}
    \vocabItem{Semicircle}{Half of a circle; area $= \frac{1}{2}\pi r^{2}$.}
\end{stepVocabBox}

\begin{stepsCard}[How to Find the Area of a Composite Shape]
    \stepItem{\textbf{Split} the figure into simple shapes you know (rectangles, triangles, circles, semicircles).}
    \stepItem{Find the \textbf{area of each} piece using the correct formula.}
    \stepItem{\textbf{Add} the areas if pieces are joined, or \textbf{subtract} if a piece is cut out.}
\end{stepsCard}

% --- Worked Example 1 ---
\begin{stepExample}{A figure is a rectangle $8 \times 4$ m with a triangle (base $8$ m, height $3$ m) on top. Find the total area.}
    \stepShow{1} Split: one rectangle and one triangle.

    \stepShow{2} Rectangle: $8 \times 4 = 32$ m$^{2}$. \quad Triangle: $\frac{1}{2} \times 8 \times 3 = 12$ m$^{2}$.

    \stepShow{3} Add: $32 + 12 = 44$ m$^{2}$.
    \stepResult{$44$ m$^{2}$}
\end{stepExample}

% --- Worked Example 2 ---
\begin{stepExample}{A rectangle $10 \times 6$ cm has a semicircle (diameter $6$ cm) attached to one short side. Find the total area. ($\pi \approx 3.14$)}
    \stepShow{1} Split: one rectangle and one semicircle (radius $= 3$ cm).

    \stepShow{2} Rectangle: $10 \times 6 = 60$ cm$^{2}$. \quad Semicircle: $\frac{1}{2} \times 3.14 \times 3^{2} = \frac{1}{2} \times 3.14 \times 9 = 14.13$ cm$^{2}$.

    \stepShow{3} Add: $60 + 14.13 = 74.13$ cm$^{2}$.
    \stepResult{$\approx 74.13$ cm$^{2}$}
\end{stepExample}

\watchOut{When a shape is \textbf{cut out}, subtract its area. When it is \textbf{attached}, add it.}

% ============================================================================
% PRACTICE
% ============================================================================
\newpage
\begin{stepPractice}{Composite Shapes Practice}
    \resetProblems

    \stepPracticeHeader[step1Color]{Identify the Pieces}
    \prob A figure is made of a square and a semicircle. Name the two simple shapes. \answerBlank[3cm]
    \answer{square and semicircle}

    \stepPracticeHeader[step2Color]{Find Each Area ($\pi \approx 3.14$)}
    \begin{multicols}{2}
    \prob Rectangle $6 \times 5$ ft. Area $=$ \answerBlank[2cm]
    \answer{$30$ ft$^{2}$}
    \prob Triangle: base $6$ ft, height $4$ ft. Area $=$ \answerBlank[2cm]
    \answer{$12$ ft$^{2}$}
    \end{multicols}

    \stepPracticeHeader[step3Color]{Combine the Areas}
    \prob A rectangle $12 \times 5$ in has a triangle (base $12$ in, height $4$ in) on top. Total area $=$ \answerBlank[2.5cm]
    \answer{$84$ in$^{2}$}

    \prob A $4 \times 4$ cm square has a circle of radius $1$ cm cut out of its center. Find the remaining area. \answerBlank[2.5cm]
    \answer{$12.86$ cm$^{2}$}

    \wordProblem{A patio is a rectangle $9 \times 6$ m with a semicircle (diameter $6$ m) at one end. What is the total area?}{m$^{2}$}
    \answer{$\approx 68.13$ m$^{2}$}
\end{stepPractice}
